\begin{theorem}[Equality of ordered edges from their endpoints]%
    \label{thm:peps_cycle_assignment_support}
    \lean{TNLean.PEPS.Edge.ofAdj_eq_iff_endpoints,
        TNLean.PEPS.Edge.ofAdj_ne_of_endpoint_coordinates}
    \leanok
    Two edges of an ordered simple graph are equal precisely when their
    unordered endpoint pairs agree. A coordinate that is constant on
    the endpoints of one edge distinguishes it from an edge whose
    endpoint has a different coordinate. These are auxiliary geometric
    identities for the accessible regular blocks of
    \cite[Section~6.3]{Schuch2010PEPS}.
\end{theorem}

\begin{proof}\leanok
    Each edge is determined by its sorted endpoint pair. Equality of the
    unordered pairs is therefore equivalent to equality of the ordered
    edges. If an endpoint coordinate differs, neither possible matching
    of the endpoint pairs can identify the two edges.
\end{proof}
